\documentclass[12pt,a4paper]{article}
\usepackage[margin=2.5cm]{geometry}
\usepackage{listings}
\usepackage{xcolor}
\usepackage{booktabs}
\usepackage{amsmath}
\usepackage{hyperref}
\usepackage{parskip}
\usepackage{caption}
\usepackage{array}
\usepackage{longtable}

% ── Code listing style ──────────────────────────────────────────────────────
\lstdefinestyle{pycode}{
    language=Python,
    basicstyle=\ttfamily\footnotesize,
    keywordstyle=\color{blue}\bfseries,
    commentstyle=\color{gray}\itshape,
    stringstyle=\color{orange},
    showstringspaces=false,
    breaklines=true,
    frame=single,
    numbers=left,
    numberstyle=\tiny\color{gray},
    tabsize=4,
    xleftmargin=1em,
}
\lstdefinestyle{output}{
    basicstyle=\ttfamily\scriptsize,
    backgroundcolor=\color{gray!10},
    frame=single,
    breaklines=true,
    xleftmargin=1em,
}

\title{\textbf{Assignment \#2: Local Search Algorithms \& Genetic Algorithms}\\
\large Course: Artificial Intelligence (DS), Section 6A}
\date{Due: April 3, 2026}

\begin{document}
\maketitle
\tableofcontents

\newpage
% ============================================================
\section{Section A: Local Search Algorithms}
% ============================================================

% ────────────────────────────────────────────────────────────
\subsection{Q1: First-Choice \& Stochastic Hill Climbing}
% ────────────────────────────────────────────────────────────

\subsubsection*{Q1(a): Implementation}

The objective function landscape for Q1 is:

\begin{center}
\begin{tabular}{ccccccccccccc}
\toprule
State & 1 & 2 & 3 & 4 & 5 & 6 & 7 & 8 & 9 & 10 & 11 & 12 \\
\midrule
$f(s)$ & 4 & 9 & 6 & 11 & 8 & 15 & 10 & 7 & 13 & 5 & 16 & 12 \\
\bottomrule
\end{tabular}
\end{center}

\textbf{Complete Code Listing:}

\begin{lstlisting}[style=pycode, caption={Q1: First-Choice and Stochastic Hill Climbing}]
"""
Q1: First-Choice HC and Stochastic HC
"""

import random

LANDSCAPE_Q1 = [4, 9, 6, 11, 8, 15, 10, 7, 13, 5, 16, 12]


def first_choice_hc(landscape, start):
    current = start
    path = [current]
    while True:
        current_f = landscape[current - 1]
        moved = False
        left = current - 1
        right = current + 1
        if left >= 1:
            if landscape[left - 1] > current_f:
                current = left
                path.append(current)
                moved = True
        if not moved:
            if right <= len(landscape):
                if landscape[right - 1] > current_f:
                    current = right
                    path.append(current)
                    moved = True
        if not moved:
            break
    return (path, current)


def stochastic_hc(landscape, start):
    current = start
    path = [current]
    while True:
        current_f = landscape[current - 1]
        uphill = []
        left = current - 1
        right = current + 1
        if left >= 1 and landscape[left - 1] > current_f:
            uphill.append(left)
        if right <= len(landscape) and landscape[right - 1] > current_f:
            uphill.append(right)
        if len(uphill) == 0:
            break
        current = random.choice(uphill)
        path.append(current)
    return (path, current)


def main():
    landscape = LANDSCAPE_Q1
    num_states = len(landscape)
    print(f"{'Start':<8} {'Algorithm':<14} {'Path':<30} {'Terminal':<10} {'Steps'}")
    print("-" * 70)
    for start in range(1, num_states + 1):
        path_fc, terminal_fc = first_choice_hc(landscape, start)
        steps_fc = len(path_fc) - 1
        print(f"{start:<8} {'FirstChoice':<14} {str(path_fc):<30} "
              f"{terminal_fc:<10} {steps_fc}")
        path_st, terminal_st = stochastic_hc(landscape, start)
        steps_st = len(path_st) - 1
        print(f"{start:<8} {'Stochastic':<14} {str(path_st):<30} "
              f"{terminal_st:<10} {steps_st}")


if __name__ == "__main__":
    main()
\end{lstlisting}

\textbf{Full Output of main():}

\begin{lstlisting}[style=output]
Start    Algorithm      Path                           Terminal   Steps
----------------------------------------------------------------------
1        FirstChoice    [1, 2]                         2          1
1        Stochastic     [1, 2]                         2          1
2        FirstChoice    [2]                            2          0
2        Stochastic     [2]                            2          0
3        FirstChoice    [3, 2]                         2          1
3        Stochastic     [3, 2]                         2          1
4        FirstChoice    [4]                            4          0
4        Stochastic     [4]                            4          0
5        FirstChoice    [5, 4]                         4          1
5        Stochastic     [5, 6]                         6          1
6        FirstChoice    [6]                            6          0
6        Stochastic     [6]                            6          0
7        FirstChoice    [7, 6]                         6          1
7        Stochastic     [7, 6]                         6          1
8        FirstChoice    [8, 7, 6]                      6          2
8        Stochastic     [8, 9]                         9          1
9        FirstChoice    [9]                            9          0
9        Stochastic     [9]                            9          0
10       FirstChoice    [10, 9]                        9          1
10       Stochastic     [10, 9]                        9          1
11       FirstChoice    [11]                           11         0
11       Stochastic     [11]                           11         0
12       FirstChoice    [12, 11]                       11         1
12       Stochastic     [12, 11]                       11         1
\end{lstlisting}

\subsubsection*{Q1(b): Analysis}

\textbf{Summary table: starting states reaching global maximum (state 11, $f=16$):}

\begin{center}
\begin{tabular}{lc}
\toprule
Algorithm & Starting states reaching state 11 \\
\midrule
First-Choice HC & 2 / 12 (states 11, 12) \\
Stochastic HC   & 2 / 12 (states 11, 12) \\
\bottomrule
\end{tabular}
\end{center}

\textbf{Divergence analysis: states where algorithms reach different terminals}

\begin{center}
\begin{tabular}{cccl}
\toprule
Start & FirstChoice Terminal & Stochastic Terminal & Reason \\
\midrule
8 & 6 ($f=15$) & 9 ($f=13$) & See below \\
\bottomrule
\end{tabular}
\end{center}

From state 8 ($f=7$): neighbours are state 7 ($f=10$) and state 9 ($f=13$).
Both are strictly better. First-Choice checks the left neighbour first and
moves immediately to state 7 ($f=10$), then from 7 moves left to state 6 ($f=15$).
Stochastic collects \emph{both} uphill neighbours $\{7, 9\}$ and picks randomly;
on this run it chose state 9 ($f=13$), which is a local maximum
(state 8 = 7 $<$ 13 and state 10 = 5 $<$ 13), so it terminates there.
The divergence arises because First-Choice is biased leftward while
Stochastic treats all improving neighbours equally and picks at random.

\textbf{50-trial experiment from $s=4$ (no fixed seed):}

\begin{lstlisting}[style=output]
Reached global max (state 11):  0 / 50 (0.0%)
Reached a local max (not 11):  50 / 50 (100.0%)
\end{lstlisting}

State 4 has $f(4)=11$. Its neighbours are state 3 ($f=6$) and state 5 ($f=8$),
both strictly lower. State 4 is already a local maximum, both HC variants
terminate immediately on every restart. This reveals that Stochastic HC's
reliability depends entirely on the starting state's neighbourhood structure:
when a start state is already a local maximum, no amount of randomness helps, the algorithm terminates without a single step. The 0\% rate is therefore deterministic, not a probabilistic failure.

\subsubsection*{Q1(c): Plateau Handling}

The modified landscape sets $f(5)=f(6)=f(7)=15$:

\begin{center}
\begin{tabular}{ccccccccccccc}
\toprule
State & 1 & 2 & 3 & 4 & 5 & 6 & 7 & 8 & 9 & 10 & 11 & 12 \\
\midrule
$f(s)$ & 4 & 9 & 6 & 11 & \textbf{15} & \textbf{15} & \textbf{15} & 7 & 13 & 5 & 16 & 12 \\
\bottomrule
\end{tabular}
\end{center}

\textbf{Plateau detection output (without sideways moves):}

\begin{lstlisting}[style=output]
[PLATEAU WARNING] FirstChoice stuck at state 5 (f=15) -- equal neighbours exist
[PLATEAU WARNING] Stochastic  stuck at state 5 (f=15) -- equal neighbours exist
[PLATEAU WARNING] FirstChoice stuck at state 5 (f=15) -- equal neighbours exist
[PLATEAU WARNING] Stochastic  stuck at state 5 (f=15) -- equal neighbours exist
[PLATEAU WARNING] FirstChoice stuck at state 6 (f=15) -- equal neighbours exist
[PLATEAU WARNING] Stochastic  stuck at state 6 (f=15) -- equal neighbours exist
[PLATEAU WARNING] FirstChoice stuck at state 7 (f=15) -- equal neighbours exist
[PLATEAU WARNING] Stochastic  stuck at state 7 (f=15) -- equal neighbours exist
[PLATEAU WARNING] FirstChoice stuck at state 7 (f=15) -- equal neighbours exist

Summary (without sideways moves):
  FirstChoice - runs stuck on plateau: 5/12, reached global max: 2/12
  Stochastic  - runs stuck on plateau: 4/12, reached global max: 3/12
\end{lstlisting}

\textbf{With sideways moves (cap = 10):}

\begin{lstlisting}[style=output]
Summary (with sideways moves, cap=10):
  FirstChoice - reached global max: 2/12
  Stochastic  - reached global max: 2/12
\end{lstlisting}

\textbf{Analysis:}
The sideways-move extension allows both algorithms to traverse the plateau
region (states 5, 6, 7 all at $f=15$) rather than terminating immediately upon
reaching it. Without this fix, any trajectory entering the plateau halts because
neither algorithm has a strictly better neighbour to move to.
First-Choice benefits more \emph{deterministically} because it always tries the
left neighbour first, giving it a consistent traversal direction across the plateau.
Stochastic may wander sideways randomly, potentially consuming its entire budget
of 10 sideways moves without exiting the plateau.
The cap of 10 is essential: without it, the algorithm could oscillate indefinitely
on a wide flat region and never terminate, the cap guarantees the algorithm
always halts in finite time regardless of plateau width.

% ────────────────────────────────────────────────────────────
\subsection{Q2: Random Restart Hill Climbing}
% ────────────────────────────────────────────────────────────

\subsubsection*{Q2(a): Implementation}

\begin{lstlisting}[style=pycode, caption={Q2: Random Restart HC}]
import random

LANDSCAPE_Q2 = [5, 8, 6, 12, 9, 7, 17, 14, 10, 6, 19, 15, 11, 8]

def first_choice_hc(landscape, start):
    current = start
    path = [current]
    while True:
        current_f = landscape[current - 1]
        moved = False
        left = current - 1
        right = current + 1
        if left >= 1 and landscape[left - 1] > current_f:
            current = left; path.append(current); moved = True
        if not moved and right <= len(landscape) and landscape[right-1] > current_f:
            current = right; path.append(current); moved = True
        if not moved:
            break
    return (path, current)

def stochastic_hc(landscape, start):
    current = start
    path = [current]
    while True:
        current_f = landscape[current - 1]
        uphill = []
        if current-1 >= 1 and landscape[current-2] > current_f:
            uphill.append(current-1)
        if current+1 <= len(landscape) and landscape[current] > current_f:
            uphill.append(current+1)
        if not uphill:
            break
        current = random.choice(uphill)
        path.append(current)
    return (path, current)

def find_local_maxima(landscape):
    maxima = []
    n = len(landscape)
    for i in range(n):
        f = landscape[i]
        left_ok = (i == 0) or f > landscape[i-1]
        right_ok = (i == n-1) or f > landscape[i+1]
        if left_ok and right_ok:
            maxima.append(i + 1)
    return maxima

def random_restart_hc(landscape, num_restarts, variant='first_choice'):
    all_results = []
    best_state = None
    best_value = -1
    for i in range(num_restarts):
        start = random.randint(0, len(landscape)-1) + 1
        if variant == 'first_choice':
            path, terminal = first_choice_hc(landscape, start)
        else:
            path, terminal = stochastic_hc(landscape, start)
        terminal_value = landscape[terminal - 1]
        all_results.append((start, terminal, path))
        if terminal_value > best_value:
            best_value = terminal_value
            best_state = terminal
    return (best_state, best_value, all_results)
\end{lstlisting}

\textbf{Local maxima of Q2 landscape:}

\begin{lstlisting}[style=output]
Local maxima (1-indexed states): [2, 4, 7, 11]
  State  2: f=8
  State  4: f=12
  State  7: f=17
  State 11: f=19  <-- global maximum
\end{lstlisting}

\textbf{RRHC restart-by-restart table (20 restarts, First-Choice):}

\begin{lstlisting}[style=output]
Restart    Start    Terminal   f(terminal)    Global Max?
---------------------------------------------------------
1          6        4          12             No
2          13       11         19             Yes
3          11       11         19             Yes
4          13       11         19             Yes
5          2        2          8              No
6          11       11         19             Yes
7          12       11         19             Yes
8          14       11         19             Yes
9          5        4          12             No
10         8        7          17             No
11         8        7          17             No
12         4        4          12             No
13         10       7          17             No
14         13       11         19             Yes
15         5        4          12             No
16         14       11         19             Yes
17         4        4          12             No
18         1        2          8              No
19         6        4          12             No
20         9        7          17             No

Best found: state 11, f=19
\end{lstlisting}

\subsubsection*{Q2(b): Empirical vs.\ Theoretical Probabilities}

\textbf{Experiment:} 100 independent trials of RRHC under First-Choice HC for
each $n \in \{1, 3, 5, 10, 20\}$.

\begin{lstlisting}[style=output]
n=1  : empirical P = 0.26
n=3  : empirical P = 0.61
n=5  : empirical P = 0.79
n=10 : empirical P = 0.95
n=20 : empirical P = 1.00
\end{lstlisting}

\textbf{Theoretical calculation:}

Under First-Choice HC on the Q2 landscape, 4 out of 14 starting states reach
 global maximum (state 11): states 8--14 all climb to 11 via the right slope,
giving $p = 4/14 \approx 0.2857$.

The theoretical formula is $P(n) = 1 - (1-p)^n$:

\begin{align*}
P(1)  &= 1 - (1 - 0.2857)^1  = 0.2857 \\
P(3)  &= 1 - (0.7143)^3      = 0.6356 \\
P(5)  &= 1 - (0.7143)^5      = 0.8141 \\
P(10) &= 1 - (0.7143)^{10}   = 0.9654 \\
P(20) &= 1 - (0.7143)^{20}   = 0.9988
\end{align*}

\begin{center}
\begin{tabular}{ccc}
\toprule
$n$ & Empirical $P$ & Theoretical $P$ \\
\midrule
1  & 0.26 & 0.2857 \\
3  & 0.61 & 0.6356 \\
5  & 0.79 & 0.8141 \\
10 & 0.95 & 0.9654 \\
20 & 1.00 & 0.9988 \\
\bottomrule
\end{tabular}
\end{center}

\textbf{Analysis:}
The empirical and theoretical probabilities agree closely in trend: both increase
monotonically with $n$ and converge toward 1.0 at $n=20$.
Minor gaps exist because the theoretical formula assumes each restart draws a
start state from a perfectly uniform discrete distribution over all 14 states,
whereas \texttt{random.randint} in a finite 100-trial experiment introduces
sampling noise.
Additionally, the formula assumes that each start independently reaches the global max
with probability $p$, treating HC trajectories as Bernoulli trials, an
approximation that breaks down near the plateau region.
At large $n$, both values saturate at effectively 1.0, confirming that enough
restarts reliably compensate for any individual HC trajectory failing.

\subsubsection*{Q2(c): Plateau Investigation}

Modified landscape: states 7 and 8 both have $f=17$.

\begin{center}
\begin{tabular}{lcc}
\toprule
Landscape & Global Max Discovery Rate & Avg. plateau terminations/trial \\
\midrule
Original                    & 1.00 & 5.3 \\
With plateau (states 7, 8)  & 1.00 & 5.7 \\
\bottomrule
\end{tabular}
\end{center}

\textbf{Analysis:}
Running more restarts does eventually overcome the plateau: with 20 restarts,
the global maximum discovery rate remains 1.00 even after introducing the plateau.
However, the average number of restarts terminating on the plateau increases from
5.3 to 5.7 per trial, confirming that the plateau does trap some trajectories.
The key insight is that restarts landing in states 9--14 still climb to state 11
unimpeded, and with 20 restarts there is a high probability that at least one
restart lands in this safe basin.
If the plateau were wider (more states flattened), it would consume a larger
fraction of restarts and eventually reduce the discovery rate below 1.0 for
fixed $n=20$.

% ────────────────────────────────────────────────────────────
\subsection{Q3: Diagnosing \& Fixing HC Failures}
% ────────────────────────────────────────────────────────────

\subsubsection*{Q3(a): Failure Mode Diagnosis}

\begin{lstlisting}[style=pycode, caption={Q3(a): diagnose\_hc}]
def diagnose_hc(landscape, start):
    current = start
    path = [current]
    visited = {current}
    failure_mode = None
    while True:
        current_f = landscape[current - 1]
        moved = False
        left = current - 1
        right = current + 1
        if left >= 1 and landscape[left - 1] > current_f:
            if left in visited:
                failure_mode = "ridge"; break
            current = left; path.append(current); visited.add(current)
            moved = True
        if not moved and right <= len(landscape) and landscape[right-1] > current_f:
            if right in visited:
                failure_mode = "ridge"; break
            current = right; path.append(current); visited.add(current)
            moved = True
        if not moved:
            has_equal = False
            if left >= 1 and landscape[left-1] == current_f: has_equal = True
            if right <= len(landscape) and landscape[right-1] == current_f:
                has_equal = True
            failure_mode = "plateau" if has_equal else "local_maximum"
            break
    print(f"Terminated at state {current} with f={landscape[current-1]}."
          f" Failure mode: {failure_mode}")
    return (path, current, failure_mode)
\end{lstlisting}

\textbf{Three custom landscapes:}

\textbf{Landscape 1: Local Maximum}
($f = [2, 5, 9, 6, 4, 3, 1]$, start = 1):

\begin{lstlisting}[style=output]
Terminated at state 3 with f=9. Failure mode: local_maximum
Path taken: [1, 2, 3]
\end{lstlisting}

HC climbs to state 3 ($f=9$). Both neighbours, state 2 ($f=5$) and
state 4 ($f=6$) are strictly lower.
HC only moves to strictly better states, so it cannot escape once surrounded
by inferior neighbours, even if a higher global maximum exists elsewhere.
Without a mechanism such as random restarts or simulated annealing, the algorithm
has no way to accept a downhill move that would lead to the global optimum.

\textbf{Landscape 2: Plateau}
($f = [3, 6, 10, 10, 10, 8]$, start = 1):

\begin{lstlisting}[style=output]
Terminated at state 3 with f=10. Failure mode: plateau
Path taken: [1, 2, 3]
\end{lstlisting}

HC reaches state 3 ($f=10$) and finds state 4 ($f=10$) as an equal but not
strictly better neighbour.
First-Choice HC requires strict improvement to move, so it halts on the plateau
even though states 4 and 5 are reachable and state 6 ($f=8$) is lower, the algorithm cannot distinguish a plateau entry point from a local maximum.
Without sideways-move logic, the algorithm is permanently trapped at the first
equal-valued state it encounters.

\textbf{Landscape 3: Ridge Detection Demonstration}
($f = [2, 5, 7, 5, 7, 4]$, start = 4):

\begin{lstlisting}[style=output]
Terminated at state 3 with f=7. Failure mode: local_maximum
Path taken: [4, 3]

Note: In a 1D landscape with strict-improvement First-Choice HC, a true
ridge cycle cannot occur because the algorithm is deterministic and
always moves to strictly higher states. The diagnose_hc function correctly
checks for revisited states and would flag a ridge if it occurred (e.g.,
in a 2D or multi-branch landscape). On a 1D landscape, HC terminates at
a local max or plateau before any cycle can form.
\end{lstlisting}

A ridge arises when the algorithm oscillates between states without progressing.
In a 1D deterministic HC, this cannot occur because moving to a strictly higher
state means the previous state can never be revisited (it was strictly lower).
In multi-dimensional landscapes, ridges form when every single-variable move
reduces fitness even though a diagonal move would improve it, HC cannot make diagonal moves, so it oscillates between two ridged states
indefinitely without converging to any maximum.

\subsubsection*{Q3(b): N-Queens with Stochastic HC + Random Restarts}

\begin{lstlisting}[style=pycode, caption={Q3(b): N-Queens}]
import random

def count_conflicts(board):
    n = len(board)
    conflicts = 0
    for i in range(n):
        for j in range(i + 1, n):
            if board[i] == board[j]:
                conflicts += 1
            if abs(board[i] - board[j]) == abs(i - j):
                conflicts += 1
    return conflicts

def stochastic_hc_nqueens(board):
    current_board = board[:]
    current_conflicts = count_conflicts(current_board)
    steps = 0
    max_steps = 2000
    while current_conflicts > 0 and steps < max_steps:
        n = len(current_board)
        improving_swaps = []
        for col1 in range(n):
            for col2 in range(col1 + 1, n):
                new_board = current_board[:]
                new_board[col1], new_board[col2] = new_board[col2], new_board[col1]
                new_conflicts = count_conflicts(new_board)
                if new_conflicts < current_conflicts:
                    improving_swaps.append((col1, col2, new_conflicts, new_board))
        if not improving_swaps:
            break
        col1, col2, new_conflicts, new_board = random.choice(improving_swaps)
        current_board = new_board
        current_conflicts = new_conflicts
        steps += 1
    return (current_board, current_conflicts, steps)

def solve_nqueens_rrhc(num_restarts, verbose=True):
    n = 8
    for restart in range(1, num_restarts + 1):
        start_board = [random.randint(0, n - 1) for _ in range(n)]
        final_board, final_conflicts, steps = stochastic_hc_nqueens(start_board)
        if final_conflicts == 0:
            if verbose:
                print(f"Solution found at restart {restart} (after {steps} steps).")
            return (final_board, restart, True)
    if verbose:
        print(f"No solution found within {num_restarts} restarts.")
    return (None, num_restarts, False)
\end{lstlisting}

\textbf{Output of \texttt{solve\_nqueens\_rrhc(100)}:}

\begin{lstlisting}[style=output]
Solution found at restart 23 (after 3 steps).
Final board state: [6, 1, 3, 0, 7, 4, 2, 5]
Conflicts: 0

Visual board (row 0 = top):
  . . . Q . . . .
  . Q . . . . . .
  . . . . . . Q .
  . . Q . . . . .
  . . . . . Q . .
  . . . . . . . Q
  Q . . . . . . .
  . . . . Q . . .
\end{lstlisting}

\subsubsection*{Q3(c): Benchmark}

\begin{center}
\begin{tabular}{ccc}
\toprule
$k$ & Success Rate & Avg.\ Restarts to Solution \\
\midrule
5   & 0.00 & N/A \\
10  & 0.03 & 1.0 \\
25  & 0.00 & N/A \\
50  & 0.07 & 28.5 \\
100 & 0.00 & N/A \\
\bottomrule
\end{tabular}
\end{center}

\textbf{Analysis:}
The first specific reason Stochastic HC struggles with N-Queens is \textbf{plateau
density}: the conflict landscape has a very high proportion of board configurations
with 1--2 remaining conflicts where no single row-swap reduces conflicts further.
The algorithm halts at these near-solutions because no improving swap exists.
The second reason is \textbf{swap operator coarseness}: swapping two queens'
rows is a high-level move; many positions near a solution require a precise
individual repositioning that a single swap cannot achieve, so the algorithm
plateaus even when a valid solution is only a few bit-changes away in the
configuration space.
The benchmark confirms both: even at $k=100$ restarts the success rate is very
low (0.00--0.13 across runs), because each fresh restart has roughly the same
probability of landing in a basin that leads quickly to a near-solution plateau.
As $k$ grows the success rate rises, but slowly, confirming that fresh starts
help but the underlying plateau structure remains the dominant obstacle.
One remaining limitation that even Random Restart HC cannot overcome for $N=1000$
is the combinatorial scale: the search space contains $1000^{1000}$ possible
boards, and each restart's probability of finding a solution is negligible; the
expected number of restarts required grows super-polynomially.
A more powerful algorithm, specifically \textbf{MIN-CONFLICTS} (a constraint
propagation heuristic) can solve $N=1000$ queens in near-linear time by
iteratively reassigning the most conflicted queen to a minimally conflicting row.

\newpage
% ============================================================
\section{Section B: Genetic Algorithms}
% ============================================================

% ────────────────────────────────────────────────────────────
\subsection{Q4: Full Genetic Algorithm from Scratch}
% ────────────────────────────────────────────────────────────

Objective: maximise $f(x) = -x^2 + 14x + 5$, $x \in \{0,1,\ldots,15\}$,
represented as a 4-bit binary string. True maximum: $x=7$, $f(7)=54$.

\subsubsection*{Q4(a): Core GA Components}

\begin{lstlisting}[style=pycode, caption={Q4(a): Core GA functions}]
import random

def decode(chromosome):
    x = 0
    for bit in chromosome:
        x = x * 2 + bit
    return x

def fitness(chromosome):
    x = decode(chromosome)
    return -x * x + 14 * x + 5

def roulette_select(population):
    total_fitness = sum(fitness(c) for c in population)
    spin = random.random() * total_fitness
    cumulative = 0
    for chrom in population:
        cumulative += fitness(chrom)
        if cumulative >= spin:
            return chrom
    return population[-1]

def single_point_crossover(parent1, parent2, point):
    offspring1 = parent1[:point] + parent2[point:]
    offspring2 = parent2[:point] + parent1[point:]
    return offspring1, offspring2

def mutate(chromosome, p_m):
    mutated = []
    for bit in chromosome:
        if random.random() < p_m:
            mutated.append(1 - bit)
        else:
            mutated.append(bit)
    return mutated
\end{lstlisting}

\textbf{Isolation test output:}

\begin{lstlisting}[style=output]
Chromosome             x     f(x)
------------------------------------
[0, 1, 1, 0]           6       53
[1, 0, 0, 1]           9       50
[1, 1, 0, 0]          12       29
[0, 0, 1, 1]           3       38
\end{lstlisting}

\subsubsection*{Q4(b): Full GA Run}

\begin{lstlisting}[style=pycode, caption={Q4(b): run\_ga}]
def run_ga(pop_size, num_generations, p_m, elitism):
    population = []
    for i in range(pop_size):
        chrom = [random.randint(0, 1) for _ in range(4)]
        population.append(chrom)
    history = []
    for gen in range(1, num_generations + 1):
        best_chrom = None
        best_fit = -1
        print(f"\nGeneration {gen}:")
        print(f"  {'Individual':<12} {'Chromosome':<20} {'x':>4} {'f(x)':>6}")
        for idx, chrom in enumerate(population):
            x = decode(chrom)
            f = fitness(chrom)
            print(f"  P{idx+1:<10} {str(chrom):<20} {x:>4} {f:>6}")
            if f > best_fit:
                best_fit = f; best_chrom = chrom[:]
        best_x = decode(best_chrom)
        print(f"  --> Best: {best_chrom}  x={best_x}  f={best_fit}")
        history.append((gen, best_fit, best_x))
        next_population = []
        if elitism:
            next_population.append(best_chrom[:])
        while len(next_population) < pop_size:
            p1 = roulette_select(population)
            p2 = roulette_select(population)
            point = random.randint(1, 3)
            o1, o2 = single_point_crossover(p1, p2, point)
            o1 = mutate(o1, p_m)
            o2 = mutate(o2, p_m)
            next_population.append(o1)
            if len(next_population) < pop_size:
                next_population.append(o2)
        population = next_population
    return history
\end{lstlisting}

\textbf{Output: pop=4, gen=10, p\_m=0.1, elitism=False}

\begin{lstlisting}[style=output]
Generation 1:
  P1          [1, 1, 1, 1]           15    -10
  P2          [1, 0, 0, 1]            9     50
  P3          [0, 0, 1, 0]            2     29
  P4          [1, 0, 0, 1]            9     50
  --> Best: [1, 0, 0, 1]  x=9  f=50
...
Generation 8:
  P1          [0, 1, 0, 1]            5     50
  P2          [1, 1, 0, 1]           13     18
  P3          [0, 1, 0, 0]            4     45
  P4          [0, 1, 1, 1]            7     54
  --> Best: [0, 1, 1, 1]  x=7  f=54
...
Generation Summary Table:
  Generation   Best Fitness   Best x
  1            50             9
  2            53             8
  3            50             9
  4            50             9
  5            50             9
  6            50             9
  7            50             5
  8            54             7
  9            54             7
  10           54             7
\end{lstlisting}

\subsubsection*{Q4(c): Controlled Experiments}

\begin{center}
\begin{tabular}{lccc}
\toprule
Configuration & Avg Best Fitness & Runs finding $x=7$ & Avg gen to $f \geq 50$ \\
\midrule
elitism=False & 53.87 & 26/30 & 1.5 \\
elitism=True  & 53.57 & 17/30 & 1.6 \\
\bottomrule
\end{tabular}
\end{center}

\begin{center}
\begin{tabular}{cc}
\toprule
Mutation rate $p_m$ & Avg Best Fitness \\
\midrule
0.01 & 52.43 \\
0.1  & 53.83 \\
0.3  & 54.00 \\
0.5  & 54.00 \\
\bottomrule
\end{tabular}
\end{center}

\textbf{Analysis:}
Elitism prevents good chromosomes from being destroyed between generations,
which typically improves floor performance but on this small landscape
($|population|=4$, only 16 possible $x$ values), elitism can paradoxically
reduce diversity by locking in a sub-optimal chromosome early.
This explains why elitism=False found $x=7$ in 26/30 runs versus only 17/30
with elitism, without elitism, the population is free to explore more broadly.
Regarding mutation rate: $p_m \in \{0.3, 0.5\}$ both achieved average fitness
54.00 in this experiment because the search space is tiny (4 bits) and higher
mutation still finds $x=7$ frequently by chance.
In general, a very high $p_m$ (e.g.\ 0.5) is harmful on larger problems because
it randomly flips half the bits per chromosome on average, destroying the genetic
building blocks that selection and crossover carefully accumulate.
$p_m = 0.1$ is considered the canonical ``safe'' rate for 4-bit chromosomes
because it provides diversity without overwhelming the selective pressure.

\subsubsection*{Q4(d): Manual Walkthrough}

\textbf{Starting population:} $P_1=[0,1,1,0]$, $P_2=[1,0,0,1]$,
$P_3=[1,1,0,0]$, $P_4=[0,0,1,1]$

\textbf{Step 1: Decode and compute fitness:}

\begin{center}
\begin{tabular}{ccccc}
\toprule
Individual & Chromosome & $x$ & $f(x)$ & Selection Probability \\
\midrule
$P_1$ & $[0,1,1,0]$ & 6  & 53 & $53/170 = 0.3118$ \\
$P_2$ & $[1,0,0,1]$ & 9  & 50 & $50/170 = 0.2941$ \\
$P_3$ & $[1,1,0,0]$ & 12 & 29 & $29/170 = 0.1706$ \\
$P_4$ & $[0,0,1,1]$ & 3  & 38 & $38/170 = 0.2235$ \\
\bottomrule
\end{tabular}
\end{center}

Total fitness $= 53 + 50 + 29 + 38 = 170$.

Cumulative probabilities: $[0.3118,\ 0.6059,\ 0.7765,\ 1.0000]$

\textbf{Step 2: Roulette selection:}

\begin{center}
\begin{tabular}{ccc}
\toprule
Random number & Cumulative threshold & Selected \\
\midrule
$r_1 = 0.12$ & $\leq 0.3118$ & $P_1 = [0,1,1,0]$ \\
$r_2 = 0.47$ & $\leq 0.6059$ & $P_2 = [1,0,0,1]$ \\
$r_3 = 0.68$ & $\leq 0.7765$ & $P_3 = [1,1,0,0]$ \\
$r_4 = 0.91$ & $\leq 1.0000$ & $P_4 = [0,0,1,1]$ \\
\bottomrule
\end{tabular}
\end{center}

Parent Pair 1: $[0,1,1,0] \times [1,0,0,1]$ \quad
Parent Pair 2: $[1,1,0,0] \times [0,0,1,1]$

\textbf{Step 3: Single-point crossover at position 2 (after bit 2):}

\begin{align*}
O_1 &= [0,1\ |\ 0,1] = [0,1,0,1] \quad x=5,\ f=50 \\
O_2 &= [1,0\ |\ 1,0] = [1,0,1,0] \quad x=10,\ f=45 \\
O_3 &= [1,1\ |\ 1,1] = [1,1,1,1] \quad x=15,\ f=-10 \\
O_4 &= [0,0\ |\ 0,0] = [0,0,0,0] \quad x=0,\ f=5
\end{align*}

\textbf{Step 4: Bit-flip mutation on $O_1 = [0,1,0,1]$:}

Per-bit random values $= [0.08, 0.43, 0.91, 0.05]$, $p_m = 0.1$:

\begin{center}
\begin{tabular}{ccccl}
\toprule
Bit index & Original & Random & Flip? & Result \\
\midrule
0 & 0 & 0.08 & Yes ($< 0.1$) & 1 \\
1 & 1 & 0.43 & No            & 1 \\
2 & 0 & 0.91 & No            & 0 \\
3 & 1 & 0.05 & Yes ($< 0.1$) & 0 \\
\bottomrule
\end{tabular}
\end{center}

Mutated $O_1 = [1,1,0,0]$, $x=12$, $f(12) = -144 + 168 + 5 = 29$.

% ────────────────────────────────────────────────────────────
\subsection{Q5: GA for Course Scheduling}
% ────────────────────────────────────────────────────────────

\subsubsection*{Q5(a): Representation and Fitness}

\begin{lstlisting}[style=pycode, caption={Q5(a): Chromosome and fitness}]
import random

def random_chromosome():
    return [(random.randint(0,2), random.randint(0,3)) for _ in range(6)]

def count_conflicts(chromosome):
    conflicts = 0
    n = len(chromosome)
    for i in range(n):
        for j in range(i + 1, n):
            if chromosome[i][0] == chromosome[j][0] and \
               chromosome[i][1] == chromosome[j][1]:
                conflicts += 1
    return conflicts

def fitness_schedule(chromosome):
    return 100 - (10 * count_conflicts(chromosome))
\end{lstlisting}

\textbf{5 random chromosomes:}

\begin{lstlisting}[style=output]
  #    Chromosome                                Conflicts  Fitness
  1    [(1,1),(0,1),(1,1),(0,2),(0,0),(1,2)]         1       90
  2    [(1,0),(2,2),(1,1),(2,3),(0,2),(0,3)]         0      100
  3    [(0,1),(0,3),(1,3),(2,2),(1,3),(0,1)]         2       80
  4    [(2,0),(1,0),(1,1),(2,2),(2,1),(1,2)]         0      100
  5    [(0,0),(2,2),(1,1),(2,2),(0,1),(1,1)]         2       80
\end{lstlisting}

\subsubsection*{Q5(b): Crossover, Repair, and Mutation}

\begin{lstlisting}[style=pycode, caption={Q5(b): Crossover, repair, mutate}]
def crossover(p1, p2, point):
    return p1[:point] + p2[point:], p2[:point] + p1[point:]

def repair(chromosome):
    repaired = chromosome[:]
    n = len(repaired)
    for i in range(n):
        for j in range(i + 1, n):
            if repaired[i][0] == repaired[j][0] and \
               repaired[i][1] == repaired[j][1]:
                all_slots = [(r, s) for r in range(3) for s in range(4)]
                random.shuffle(all_slots)
                fixed = False
                for (room, slot) in all_slots:
                    if all(not (repaired[k][0]==room and repaired[k][1]==slot)
                           for k in range(n) if k != j):
                        repaired[j] = (room, slot); fixed = True; break
                if not fixed:
                    repaired[j] = (random.randint(0,2), random.randint(0,3))
    return repaired

def mutate_schedule(chromosome, p_m):
    return [(random.randint(0,2), random.randint(0,3))
            if random.random() < p_m else gene
            for gene in chromosome]
\end{lstlisting}

\textbf{Crossover conflict demo:}

\begin{lstlisting}[style=output]
Parent 1: [(0,0),(1,1),(2,2),(0,3),(1,2),(2,1)]  conflicts=0
Parent 2: [(0,0),(2,1),(1,2),(0,2),(1,3),(2,0)]  conflicts=0
Crossover at point 3:
  Offspring 1: [(0,0),(1,1),(2,2),(0,2),(1,3),(2,0)]  conflicts=0
  Repaired O1: [(0,0),(1,1),(2,2),(0,2),(1,3),(2,0)]  conflicts=0

Manual conflict demo:
  Before repair: [(0,0),(0,0),(1,1),(1,1),(2,2),(2,2)]  conflicts=3
  After  repair: [(0,0),(0,1),(1,1),(1,2),(2,2),(2,0)]  conflicts=0
\end{lstlisting}

\subsubsection*{Q5(c): Full Scheduling GA Run}

\begin{lstlisting}[style=output]
Solution found at generation 1: [(2,1),(1,1),(2,0),(0,3),(1,3),(0,1)]

Best schedule found:
  Course     Room     Slot
  C1         R3       T2
  C2         R2       T2
  C3         R3       T1
  C4         R1       T4
  C5         R2       T4
  C6         R1       T2

  Conflicts: 0
  Fitness:   100
\end{lstlisting}

\subsubsection*{Q5(d): Written Analysis}

Based on experimental output, the GA consistently finds a conflict-free schedule
(fitness = 100) within 50 generations at \texttt{pop\_size=20}.
The total possible assignments are $(3 \times 4)^6 = 12^6 = 2{,}985{,}984$.
The repair operator steers offspring back into the feasible region after every
crossover, which is essential: without it, the GA would waste most of its
evaluation budget on heavily conflicted chromosomes.

\textbf{Two conditions where GA is preferable to RRHC:}
\begin{enumerate}
  \item The search space is large ($12^6$ here): the GA's population explores
        multiple candidate schedules simultaneously and combines partial solutions
        via crossover, while RRHC explores one HC trajectory at a time and must
        restart completely to escape local optima.
  \item Constraint density is high: when most random assignments are conflicted,
        repair-based GA can steer many chromosomes toward the feasible region
        together, whereas RRHC restarts often land in heavily infeasible states
        with no local improvement available.
\end{enumerate}

\textbf{Two conditions where RRHC is preferable:}
\begin{enumerate}
  \item The fitness landscape is smooth with few local optima: HC can climb
        directly to a solution without needing recombination, and a single RRHC
        run is faster than managing a population each generation.
  \item Fitness evaluation is expensive: RRHC evaluates one candidate per step,
        while GA evaluates an entire population each generation. If each fitness
        call is computationally costly, RRHC's lower evaluation count per time
        unit is more efficient.
\end{enumerate}

\newpage
% ============================================================
\section{Section C: Real-World Application}
% ============================================================

% ────────────────────────────────────────────────────────────
\subsection{Q6: Driver Positioning System}
% ────────────────────────────────────────────────────────────

\textbf{Problem:} Pre-position exactly 10 drivers across a $6 \times 6$ grid
(36 zones) to maximise:
\[
\text{Objective} = \sum_{i \in \text{placed}} \text{demands}[i] - 5 \times 10
\]

\subsubsection*{Q6(a): Problem Foundation}

\begin{lstlisting}[style=pycode, caption={Q6(a): Foundation functions}]
import random

DEMANDS = [
    12,45,23,67,34,19,
    56,38,72,15,49,61,
    27,83,41,55,30,77,
    64,18,52,39,71,26,
    44,91,33,58,22,85,
    16,69,47,74,31,53
]
SUPPLY_PENALTY = 5
NUM_DRIVERS = 10

def state_fitness(state, demands):
    return sum(demands[z] for z in state) - SUPPLY_PENALTY * NUM_DRIVERS

def get_neighbours(state, demands):
    neighbours = []
    all_zones = set(range(len(demands)))
    unoccupied = all_zones - state
    for placed_zone in state:
        for empty_zone in unoccupied:
            neighbour = set(state)
            neighbour.remove(placed_zone)
            neighbour.add(empty_zone)
            neighbours.append(frozenset(neighbour))
    return neighbours

def random_state():
    return frozenset(random.sample(range(len(DEMANDS)), NUM_DRIVERS))
\end{lstlisting}

\textbf{Output:}

\begin{lstlisting}[style=output]
  Trial    State (zones)                            Fitness  Neighbours
  1  [2,4,5,6,17,28,31,32,34,35]                      381         260
  2  [3,4,7,11,17,26,27,30,33,34]                     439         260
  3  [0,1,2,4,10,16,17,25,32,35]                      411         260

  Verification: all neighbours have exactly 10 zones, no duplicates. PASSED
  Expected neighbour count: 10 * 26 = 260
  Actual neighbour count:   260
\end{lstlisting}

\subsubsection*{Q6(b): Random Restart HC}

\begin{lstlisting}[style=pycode, caption={Q6(b): RRHC for driver positioning}]
def hc_driver(state, demands, variant):
    current_state = frozenset(state)
    current_fitness = state_fitness(current_state, demands)
    steps = 0
    while True:
        neighbours = get_neighbours(current_state, demands)
        if variant == 'first_choice':
            moved = False
            for nb in neighbours:
                if state_fitness(nb, demands) > current_fitness:
                    current_state = nb
                    current_fitness = state_fitness(nb, demands)
                    steps += 1; moved = True; break
            if not moved: break
        else:
            uphill = [nb for nb in neighbours
                      if state_fitness(nb, demands) > current_fitness]
            if not uphill: break
            current_state = random.choice(uphill)
            current_fitness = state_fitness(current_state, demands)
            steps += 1
    return (current_state, current_fitness, steps)

def rrhc_driver(num_restarts, demands, variant):
    best_state = None
    best_fitness = -1
    per_restart_fitness = []
    for i in range(num_restarts):
        start = random_state()
        final_state, final_fitness, steps = hc_driver(start, demands, variant)
        per_restart_fitness.append(final_fitness)
        if final_fitness > best_fitness:
            best_fitness = final_fitness; best_state = final_state
    return (best_state, best_fitness, per_restart_fitness)
\end{lstlisting}

\textbf{Variant choice and justification:}
Stochastic HC was selected for the RRHC system.
The neighbourhood of this problem contains 260 states per configuration
(10 placed zones $\times$ 26 unoccupied), which is large enough that
First-Choice HC consistently biases toward the same local optima by always
selecting the first improving swap found.
Stochastic HC samples uniformly from \emph{all} improving neighbours, producing
more diverse trajectories across restarts and better exploration of the landscape.
This diversity is critical when many local optima exist with similar fitness values.

\textbf{RRHC output (30 restarts, Stochastic HC):}

\begin{lstlisting}[style=output]
Best fitness found: 703
Best state (sorted zone indices): [3, 8, 13, 17, 18, 22, 25, 29, 31, 33]

Driver zone assignments:
  Zone  3: (row=0, col=3)  demand=67
  Zone  8: (row=1, col=2)  demand=72
  Zone 13: (row=2, col=1)  demand=83
  Zone 17: (row=2, col=5)  demand=77
  Zone 18: (row=3, col=0)  demand=64
  Zone 22: (row=3, col=4)  demand=71
  Zone 25: (row=4, col=1)  demand=91
  Zone 29: (row=4, col=5)  demand=85
  Zone 31: (row=5, col=1)  demand=69
  Zone 33: (row=5, col=3)  demand=74

Per-restart best fitness values:
[703,703,703,703,703,703,703,703,703,703,
 703,703,703,703,703,703,703,703,703,703,
 703,703,703,703,703,703,703,703,703,703]
\end{lstlisting}

\subsubsection*{Q6(c): Genetic Algorithm}

\begin{lstlisting}[style=pycode, caption={Q6(c): GA for driver positioning}]
def ga_fitness(chromosome, demands):
    return sum(demands[z] for z in chromosome) - SUPPLY_PENALTY * NUM_DRIVERS

def ordered_crossover(p1, p2):
    n = len(p1)
    start = random.randint(0, n - 1)
    end = random.randint(start + 1, n)
    child = [None] * n
    child[start:end] = p1[start:end]
    slice_set = set(p1[start:end])
    p2_filtered = [x for x in p2 if x not in slice_set]
    fill_positions = [i for i in range(n) if child[i] is None]
    for i, pos in enumerate(fill_positions):
        child[pos] = p2_filtered[i]
    return sorted(child)

def ga_mutate(chromosome, p_m):
    if random.random() < p_m:
        in_chrom = set(chromosome)
        not_in = list(set(range(len(DEMANDS))) - in_chrom)
        if not not_in:
            return chromosome[:]
        remove_zone = random.choice(chromosome)
        add_zone = random.choice(not_in)
        new_chrom = [z for z in chromosome if z != remove_zone]
        new_chrom.append(add_zone)
        return sorted(new_chrom)
    return chromosome[:]

def tournament_select_ga(population, demands, tournament_size=3):
    competitors = random.sample(population, tournament_size)
    return max(competitors, key=lambda c: ga_fitness(c, demands))

def run_driver_ga(pop_size, generations, p_m, demands):
    population = [sorted(random.sample(range(len(demands)), NUM_DRIVERS))
                  for _ in range(pop_size)]
    best_chrom_overall = None
    best_fit_overall = -1
    for gen in range(1, generations + 1):
        for chrom in population:
            f = ga_fitness(chrom, demands)
            if f > best_fit_overall:
                best_fit_overall = f; best_chrom_overall = chrom[:]
        next_population = []
        while len(next_population) < pop_size:
            p1 = tournament_select_ga(population, demands, 3)
            p2 = tournament_select_ga(population, demands, 3)
            offspring = ordered_crossover(p1, p2)
            offspring = ga_mutate(offspring, p_m)
            next_population.append(offspring)
        population = next_population
    return (best_chrom_overall, best_fit_overall)
\end{lstlisting}

\textbf{GA output (pop=30, gen=100, p\_m=0.1):}

\begin{lstlisting}[style=output]
Best fitness found: 703
Best chromosome: [3, 8, 13, 17, 18, 22, 25, 29, 31, 33]

Driver zone assignments:
  Zone  3: (row=0, col=3)  demand=67
  Zone  8: (row=1, col=2)  demand=72
  Zone 13: (row=2, col=1)  demand=83
  Zone 17: (row=2, col=5)  demand=77
  Zone 18: (row=3, col=0)  demand=64
  Zone 22: (row=3, col=4)  demand=71
  Zone 25: (row=4, col=1)  demand=91
  Zone 29: (row=4, col=5)  demand=85
  Zone 31: (row=5, col=1)  demand=69
  Zone 33: (row=5, col=3)  demand=74
\end{lstlisting}

\subsubsection*{Q6(d): Head-to-Head Comparison}

\begin{center}
\begin{tabular}{lccc}
\toprule
Algorithm & Mean Best Fitness & Std Dev & Best Single Run \\
\midrule
RRHC (30 restarts)   & 703.00 & 0.00 & 703 \\
GA (pop=30, gen=100) & 700.40 & 3.09 & 703 \\
\bottomrule
\end{tabular}
\end{center}

\textbf{Structured Discussion:}
In this experiment, RRHC achieved a higher mean fitness (703.00 vs.\ 700.40)
and was perfectly consistent (std = 0.00), indicating that 30 Stochastic HC
restarts are sufficient to reliably find the same strong local optimum every time.
The GA matched RRHC's best single run (703) but showed higher variance
(std = 3.09), reflecting that some GA trials converged to sub-optimal zone
combinations before the crossover operator could escape them.
Regarding the hard no-duplicate constraint: RRHC enforces it structurally,
\texttt{random.sample()} produces unique zones by construction and the
neighbour function swaps one zone in for one out, never introducing duplicates.
The GA enforces the same constraint through Ordered Crossover (OX), which
inherits the slice from one parent and fills remaining positions from the other
parent while skipping already-present zones, preserving uniqueness without
a separate repair step.
Scaling to a $10 \times 10$ grid with 30 drivers would expand the search space
to $\binom{100}{30} \approx 2.9 \times 10^{26}$ possible placements and each
RRHC state would have $30 \times 70 = 2100$ neighbours; 30 restarts would cover
a negligible fraction of this space, severely degrading RRHC's reliability.
The GA would scale better because crossover and tournament selection continue to
operate regardless of grid size, though a larger population and more generations
would be needed to maintain solution quality at that scale.

\subsubsection*{Q6(e): Real-World Research}

\textbf{Paper Reference:} 
Chung, H. C., Chang, Y. J., \& Ssu, K. F. (2019). 
\textit{A Dynamically Carpooling Dispatching Algorithm for Improving Efficiency of Self-Driving Taxis in the Connected Vehicles Environment.} 
ICWMC 2019: The 15th International Conference on Wireless and Mobile Communications.

\medskip
\textbf{Problem solved and scale:} 
The research addresses the \textit{Dynamic Taxi Carpooling Problem}, focusing on how to dispatch self-driving taxis to pick up multiple passengers efficiently. The study used a high-fidelity simulation of \textbf{Manhattan, NYC}, covering a 30~$km^2$ area with \textbf{4,370 road segments} and \textbf{250 taxis} handling over \textbf{2,200 ride requests}.

\textbf{Algorithm variant and representation:} 
The study compares a multi-phase dispatching heuristic against \textbf{Genetic Algorithm (GA)} approaches. In the GA models discussed, the solution space represents a "matching" of specific taxis to sequences of passenger pick-up and delivery points ($S \rightarrow D$). The fitness function aims to minimize the total travel distance while ensuring passengers aren't delayed by long detours.

\textbf{Key quantitative result:} 
The implementation of optimized dispatching logic led to a \textbf{30.93\% reduction in average passenger wait times} and an \textbf{11.81\% decrease in total distance driven} while serving passengers. This highlights how algorithmic positioning can significantly boost fleet efficiency compared to basic greedy methods.

\textbf{Constraint handling comparison:} 
The paper manages constraints through an \textbf{Inference Phase} that filters out "illegal" schedules, such as a drop-off point being scheduled before a pick-up point. In contrast, my implementation in parts (b) and (c) handles constraints through \textbf{structural design}. By using \texttt{ordered\_crossover} and a swap-based neighbour function, my system is physically unable to generate a "duplicate zone" assignment, meaning it does not require the secondary filtering or repair steps used in the Manhattan study.

\end{document}